\odiseachapter{La Gran Guerra}{La Gran Guerra}\label{cap:guerra}

Permíteme, tolerante lectora, generoso lector, que en este capítulo abandone a Homero para seguir las cavilaciones del anciano Odiseo de \emph{Abandonando Ítaca}. Imagínatelo sentado junto al fuego, enjuto, su antigua corpulencia convertida en tendón y nervio, los ojos todavía feroces a pesar del temblor en las manos y el andar encorvado, la mente todavía lúcida, quizás más lúcida que nunca. Son las cuatro de la madrugada y los recuerdos, o los fantasmas, no le dejan dormir. Se levanta en silencio y se sienta frente al fuego, con una copa de vino tibio entre las manos. Recuerda aquellos días, cuando contó al rey Alcínoo cómo empezó la gran guerra.

%
\begin{verse}
Durante muchas noches has contado tu historia,\\
en la corte del padre de Nausícaa.\\
Era un rey, pero tu memoria se niega a recordar el nombre.\\
Te pasa cada vez más a menudo, primero le pasó a Laertes,\\
olvidas el nombre de las personas,\\
ya no puedes recordar las caras,\\
ni puedes decir lo que hiciste ayer.\\
En cambio, el pasado se muestra con colores tan vivos,\\
que a menudo temes que los muertos se adueñen de los vivos.
\stanzabreak
Sí, les contaste cómo empezó la Gran Guerra.\\
Las razones fueron intrascendentes,\\
y sin embargo tan inevitables.\\
¿Cómo no creer en el destino, recordando cómo empezó todo?\\
¿Cómo no pensar que todos los hombres son juguetes de la fatalidad?
\end{verse}
%

Odiseo recuerda el juramento que los príncipes aqueos habían hecho y que les obligaba a defenderse los unos a los otros, el juramento que invoca Agamenón para convocarlos:
%
\begin{verse}
Ninguno de vosotros, reyes, se preocupó por Menelao,\\
---que era más dulce con sus caballos que con su esposa---,\\
y aún menos os gustaba su hermano, el retorcido Agamenón.\\
Pero todos estabais atados por un juramento o, al menos,\\
eso dijisteis a vuestros hijos y vuestras esposas\\
y eso declarasteis en los templos, mientras vitoreaban vuestros súbditos\\
y los sacerdotes asentían con solemnidad.
\stanzabreak
El juramento era real. En la boda de Helena,\\
se pidió a los pretendientes que establecieran un vínculo irrompible,\\
con el que triunfase, ayudándole si alguna vez surgiese la necesidad.\\
Y todos os comprometisteis, poniendo por testigos a los dioses.\\
¿Cómo podíais no hacerlo? Todos estabais convencidos\\
de ser el elegido. Y el trato era dulce para el ganador, desposar\\
la belleza y ganar el apoyo de los otros príncipes,\\
¿quién podría resistirse a un trato así? No tú, orgulloso Ulises,\\
el más inteligente de todos vosotros.
\end{verse}
%
Viejo y amargado, Ulises no ha perdido la lucidez ni la capacidad de reírse de sí mismo.
%
\begin{verse}
Últimamente has llegado a preguntarte\\
si eras, después de todo, tan inteligente.\\
Curioso, ahora que tus pensamientos\\
son a menudo confusos y necesitas un esfuerzo para conversar,\\
ahora que te adormeces tanto y que tu espíritu deambula con frecuencia,\\
ahora que ves en los ojos de la familia y de los sirvientes,\\
la compasión ---¿o es la condescendencia?--- hacia el anciano,\\
apenas disfrazada de respetuosa preocupación,\\
has llegado a dudar de muchas cosas\\
que eran tan obvias en tu juventud.
\stanzabreak
Y una de ellas es tu supuesta astucia.\\
Claro que eras artero, y avispado,\\
mucho más que los otros.\\
Pero ¿era un logro tan grande?\\
Al fin y al cabo, la mayoría de ellos eran montañas de carne,\\
como el bueno de Áyax, o incluso Aquiles el divino,\\
todo músculos y sin ningún ingenio. Pero, si eras tan inteligente,\\
¿cómo no entendiste que la más hermosa mujer\\
que jamás pisó la Tierra no iba a ser otorgada\\
al pobre jefe de un islote, que solo podía ofrecer\\
cabras, bravuconadas y sus trucos de sabelotodo?
\stanzabreak
Finalmente, Helena fue para el mejor\\
postor, como ha ocurrido siempre.\\
Quizás se debatiese si era una mejor opción\\
Aquiles, pero la verdad es que un gran héroe\\
no es tan buena inversión como un rey poderoso.\\
Así, Helena desposó al espartano,\\
y todos regresaron a casa sintiéndose engañados,\\
---incluso tú, que te habías forjado fama de embustero---\\
y sin ninguna intención de cumplir vuestras promesas.
\end{verse}

La competición por la mano de la bella Helena y el pacto posterior de los reyes no aparecen en Homero, sino bastante más tarde, en el \emph{Catálogo de las mujeres} atribuido a Hesíodo (siglo VI a.C.), pero, dado cómo se las gastaban los caudillos aqueos, podía haber salido de la pluma ---quiero decir el punzón--- del Poeta. La vuelta de tuerca que nos ofrece el anciano Odiseo es esta:
%

\begin{verse}
Solo que, cuando Agamenón llamó,\\
no solamente recordó el juramento\\
e invocó vuestro honor. No, también describió\\
una maravillosa Troya, la ciudad de altos muros\\
e incontables riquezas. Así que ayudar al amigo traicionado\\
vino con un buen extra. Después de todo, el gran ejército aqueo\\
no debería tener problemas para asaltar el lugar en pocas semanas,\\
y los líderes se repartirían el inimaginable botín.
\stanzabreak
Todos se engancharon con facilidad,\\
al ser la codicia más poderosa que la envidia y el rencor,\\
y, además, a todos os alegraba que Helena\\
hubiese hecho cornudo al soberbio espartano;\\
también os hacía felices vengar la ofensa a sangre y fuego,\\
para que vuestras esposas, vuestras hijas y hermanas,\\
tomasen buena nota de quién estaba al mando.
\end{verse}

Este Ulises tiene muy poco que ver con el santo varón que nos pinta Nolan, obligado a seguir a Agamenón para evitar represalias. Pero sigamos con la guerra, cuyas imágenes parece conjurar el crepitar del fuego que apenas calienta los cansados huesos del anciano.

\begin{verse}
¿Qué esperabas?\\
¿Realmente pensaste que podías tomar Troya,\\
con sus altos muros y valientes soldados,\\
con su orgulloso rey y ese héroe gigante,\\
el príncipe Héctor, domador de caballos,\\
mostrándote y pidiéndoles su rendición?
\stanzabreak
La verdad es\\
que casi se burlaron de vosotros.\\
Eran tan fuertes y tan poderosos\\
como para no temer a esa banda de bárbaros\\
que pululaban más allá de sus altos muros cada año.
\end{verse}

Un pequeño error de cálculo, no comprender que Troya era, en realidad, inexpugnable. Un pequeño error de cálculo que les habría costado la derrota, de hecho, si no hubiera sido por el capricho de los dioses. Pero el anciano Odiseo no cree en los dioses y está convencido de que los mortales se bastan y se sobran para buscarse la ruina.

\begin{verse}
Quizá, pensándolo mejor,\\
la perdición de Troya no fue debida a vuestro truco,\\
sino a su orgullo. Nunca consideraron, ni por un momento,\\
que su enemigo tuviese alguna posibilidad.\\
Por eso os toleraron durante diez largos años,\\
en lugar de golpear a vuestro ejército con todas sus fuerzas,\\
incendiando vuestros barcos\\
y sembrando de cadáveres el mar.
\stanzabreak
Y, sin embargo, el viejo Príamo era, a su manera,\\
tan artero como Agamenón,\\
reacio a sacrificar a demasiados hombres\\
para barreros de sus playas para siempre.\\
Después de todo, él pensó\\
que derrotar al gran Aquiles costaría demasiada sangre,\\
así que ¿por qué no hacer que muriesen de hambre?,\\
¿por qué no debilitar a los bárbaros con pequeñas escaramuzas\\
y dejar que el sol de Apolo y la peste hiciesen su trabajo?
\stanzabreak
Sí, demasiado soberbios.\\
¿Quién en la Tierra sería tan estúpido\\
como para aceptar un regalo de enemigos jurados\\
y llevarlo detrás de las murallas?
\end{verse}

Un pequeño error de cálculo que les costaría una década baldía:

\begin{verse}
Transcurrieron así diez largos años.\\
En sus tiendas, los reyes aqueos\\
enviaban hombres a morir cada mañana\\
y ordenar cada noche enterrar los cadáveres,\\
mientras se repetían las consignas\\
que ya nadie creía.
\stanzabreak
En su palacio, el viejo Príamo\\
finge que los aqueos no existen o que son\\
simplemente un deporte para sus muchos hijos,\\
sin hacer caso a Héctor,\\
al parecer el único\\
que ve el peligro.
\stanzabreak
En su lecho, después de tanto tiempo,\\
Helena se pregunta qué es lo que vio\\
en ese Paris debilucho,\\
aunque era bello, sí, mas tan débil y vano,\\
tan charlatán, tan insignificante.
\stanzabreak
De noche, Menelao, el de los anchos hombros,\\
echa de menos más a sus caballos que a su esposa,\\
y Agamenón intenta olvidar su culpa,\\
ahogándola en pecados cada vez más grandes.\\
Ahora se acuesta con Criseida,\\
sabiendo muy bien que está prohibida por Apolo\\
y que puede traer problemas innecesarios a su aventura.\\
Incluso Aquiles, que disfruta matando cada día laborable\\
su cuota de troyanos, se pregunta si quizás ya es hora\\
de sacar a sus soldados y a su joven amante de la aburrida guerra\\
y regresar a casa.
\stanzabreak
En cuanto a ti, astuto Ulises,\\
¿qué hacías además de masacrar hombres y beber vino?\\
Tu esposa estaba sola,\\
tu hijo crecía como un huérfano\\
y tú desperdiciabas tus años dorados,\\
sin poder escapar de la atracción de la llama,\\
como una polilla que vuela dando vueltas en torno a la vela.
\end{verse}

Diez años que podrían haber sido veinte. O quizás estabais ya a punto de renunciar y volver a casa, cuando Agamenón decidió complicar las cosas más de lo que ya estaban.


\begin{verse}
Y entonces, después de diez años baldíos,\\
Agamenón no tuvo mejor idea que presumir\\
en el peor momento. Ah, pero no pudo resistir la tentación\\
de humillar al gran héroe, que, después de todo,\\
parecía ser el único que se divertía en el conflicto.
\stanzabreak
Por supuesto, Aquiles y sus mirmidones se unían a cada batalla\\
y mataban muchos más troyanos que cualquiera,\\
pero, al mismo tiempo, era evidentemente\\
solo un deporte para ellos. El semidiós debía cumplir su destino\\
y quedó tan abandonado como todos los demás en las costas de Ilion,\\
pero, a diferencia de los demás, no estaba interesado en el futuro.
\stanzabreak
¿Por qué habría de estarlo? La profecía era clara\\
como las aguas de Xanto. Su vida iba a ser corta y gloriosa\\
y la gloria la tenía cada día\\
---aplastando huesos y desgarrando vientres,\\
atravesando corazones y reventando cráneos,\\
haciendo centenares de viudas y huérfanos---\\
mientras esperaba su destino.\\
No tenía anhelos ni remordimientos,\\
todo lo que le era valioso estaba al alcance de su mano,\\
sus hombres, sus seres queridos y esa joven, Briseida,\\
que no debería haber significado nada para él ---una simple esclava---,\\
pero a quien amaba casi tanto como a Patroclo.
\end{verse}

Quizás fue la paz interior de Aquiles la que desquició al caudillo, incapaz de conciliar el sueño mientras veía desmoronarse sus ambiciones:

\begin{verse}
Así que Aquiles estaba contento y ese retorcido rey\\
se resintió por no ganar la guerra de la noche a la mañana,\\
como todos vosotros esperabais. Por desgracia, no podía\\
dejar pasar la oportunidad de herir al héroe\\
mientras se procuraba un nuevo juguete para sus noches.
\stanzabreak
Tú, que no podrías decir por todo el oro del mundo\\
lo que ha entrado en el plato de tu desayuno esta mañana,\\
recuerdas, sin embargo, con nitidez la escena.\\
El insolente micénico, lleno de sí mismo, sacando pecho,\\
desafiando al hijo de Peleo delante de todos los aqueos. ¿Para qué?\\
¿No entendía el tonto las implicaciones?\\
Tal vez sí, pero el orgullo y la lujuria\\
siempre se han salido con la suya en hombres poderosos.\\
Y así Briseida fue a la tienda de Agamenón\\
y Aquiles dio por terminada su guerra personal.
\end{verse}

El resto de la historia, piensa Odiseo, la repiten los bardos, convenientemente endulzada. Uno de ellos ha pasado esa misma noche por su palacio y cantado los versos que ya repite todo el mundo\footnote{Naturalmente se trata de los versos iniciales de la \emph{Ilíada}. La traducción es mía.}:
\begin{verse}
Cántame, diosa, la horrible rabia\\
del gran Aquiles, hijo de Peleo,\\
que tanto dolor causó a los argivos\\
y envió de nobles héroes las almas\\
al Hades, sus cuerpos de perros presa\\
y de pájaros banquete, cumpliendo\\
el artero plan de Zeus juntanubes,\\
enfrentó a Agamenón, de hombres líder,\\
y al glorioso Aquiles de pies ligeros.
\end{verse}
